\chapter{The Shape of Persistence}

\nopagebreak
\epigraph{"The limiting factor on statistics projects has rarely been my knowledge of statistics. The limiting factors are things like communication, endurance, organization."}{John D. Cook}
\nopagebreak

By 2025, the numbers had reached 5,233 posts and 1.8 million words, with influence that extends beyond measurement. The full shape of the archive reveals what made everything else possible: persistence.

\bigskip

\section*{The Pandemic Invariance}

March 2020 upended most patterns. Schools closed. Offices emptied. Conferences cancelled. World rhythms broke.

Blog rhythms did not.

\begin{longtable}{p{3.00cm}p{3.00cm}p{3.00cm}}
\toprule
\textbf{Year} & \textbf{Posts} & \textbf{Avg Words} \tabularnewline
\midrule
2019 & 288 & 446 \tabularnewline
2020 & 288 & 451 \tabularnewline
2021 & 341 & 451 \tabularnewline
\bottomrule
\end{longtable}

2020 is indistinguishable from the surrounding years. Twenty-two posts in March (the month everything changed), 21 in April, 23 in May. No spike of anxiety-posting, no drought of distraction. Just another year.

The flatness is striking. The practice serves a deeper function: a place where external chaos gets ordered into mathematical logic.

Cook rarely commented on the pandemic directly. A handful of posts used COVID-19 data as statistical examples. Otherwise, he simply continued. Another post about probability. Another about Python. Another about number theory.

The world outside changed, but the blog's function did not.

\bigskip

\section*{No Gap Longer Than Eleven Days}

Every year, every season, for nearly eighteen years.

Longest gap: July 9-21, 2014 (11 days).
Second-longest gaps: 8 days, appearing three times.

No gap reaches two weeks. Most professionals take longer vacations. A two-week holiday would set a record.

This isn't workaholism---the posting frequency is measured, not manic. It's something else: a commitment to showing up regardless of mood or circumstance.

Whatever else happened---travel, illness, family, career changes---the blog kept appearing.

\bigskip

\section*{The Trade-Off That Didn't Hold}

Early data suggested a trade-off: depth versus frequency. Fewer posts meant longer posts.

Later data broke the pattern:

\begin{longtable}{p{1.80cm}p{1.80cm}p{1.80cm}p{1.80cm}p{1.80cm}}
\toprule
\textbf{Era} & \textbf{Years} & \textbf{Posts} & \textbf{Avg Words} & \textbf{Total Words} \tabularnewline
\midrule
Early & 2008-2012 & 1,734 & 232 & 401,870 \tabularnewline
Middle & 2013-2017 & 968 & 297 & 287,865 \tabularnewline
Recent & 2018-2025 & 2,531 & 448 & 1,133,465 \tabularnewline
\bottomrule
\end{longtable}

The recent era produced more posts AND longer posts---more total words than the first two eras combined. The assumed trade-off didn't hold.

This is what mastery looks like: early years learning the form, middle years consolidating, recent years producing more with more substance and without apparent strain.

\bigskip

\section*{Reader Encounters}

The influence lives in encounters. A sample:

\textbf{tom\_b (Hacker News):}
\begin{blogquote}
"I have read John Cook's blog for some time and think he might be a good fit as a consultant... I suspect it would be worth your time to get a consultant like him to come in for a week."
\end{blogquote}

\textbf{RcouF1uZ4gsC (Hacker News):}
\begin{blogquote}
"If I was looking for a consultant in these matters, he would be at the top of my list."
\end{blogquote}

\textbf{wenc (Hacker News):}
\begin{blogquote}
"John D Cook's business is doing well... built reputation through insightful blog posts over years, establishing credibility in the applied mathematics community."
\end{blogquote}

\textbf{dustingetz (Dustin Getz, founder of Hyperfiddle):}
\begin{blogquote}
"Best two advanced software engineering blogs ever."
\end{blogquote}

\textbf{cperciva (Colin Percival, creator of Tarsnap and scrypt, FreeBSD release engineer):}
\begin{blogquote}
"If all three measurements of X are smaller than all three measurements of Y, you have X < Y with 95\% confidence."
\end{blogquote}
\textit{(Applying Cook's "Rule of Three" to performance testing)}


\textbf{an1sotropy (Hacker News):}
\begin{blogquote}
"I first learned about 'leaky abstractions' from John Cook, who describes IEEE 754 floats as a leaky abstraction of the reals."
\end{blogquote}

The shared theme across these comments is trust---built through consistency, through showing up day after day with useful explanations.

\bigskip

\section*{The Full Arc}

From 2008 to 2025:

\textbf{2008-2012 (The Notebook):} 1,734 posts averaging 232 words. A consultant's scratchpad. Scattered topics. Short observations. Establishing the habit.

\textbf{2013-2017 (The Crucible):} 968 posts averaging 297 words. The quiet year. The career transition. The shedding of professional obligations. Finding the mathematical core.

\textbf{2018-2025 (The Institution):} 2,531 posts averaging 448 words. Focused. Deep explanations. Interconnected references. Infrastructure.

The trajectory wasn't linear---it dipped in 2014 and transformed after 2016---but the direction was clear: from scattered to focused, from brief to substantial, from notebook to institution.

\bigskip

\section*{What the Numbers Show}

5,233 posts.
1,823,200 words.
6,533 days.
348 words per post (mean).
312 words per post (median).

The mean exceeds the median by 36 words---a gap revealing the long tail of deep-dive essays pulling the average up. The blog is not uniformly brief; a core of substantial posts defines its value.

Tuesday leads overall. Morning dominates consistently. The rhythm persists even as content changes.

\bigskip

\section*{For Anyone Starting}

What sustained practice produces:

\textbf{Start before you're ready.} The 2008 blog wasn't polished---it was exploratory. The polish came later.

\textbf{Let interests compete.} Early breadth wasn't a mistake; it was an experiment. The middle-era narrowing happened because the experiment produced a winner.

\textbf{Trust the process to clarify.} The 2014 dip wasn't failure---it was focus emerging. Sometimes fewer posts means better posts.

\textbf{Keep showing up.} Later quality came from repetition. Explaining the gamma function and elliptic curves so many times that explanations became effortless.

\bigskip

\section*{The Question}

The archive proves the thesis: you can just do things. Write about what interests you, publish without permission, show up day after day, and let the work find its form. No plan needed, no platform needed---just time and consistency.

Not just what Cook will do next, but what anyone can build with seventeen years of consistent effort.

\bigskip

\textit{The shape is complete. The arc is traced. But some things resist quantification. The final chapter acknowledges what counting misses.}

